\epigraph{"The axioms for a groups are short and natural... Yet somehow hidden behind these axioms is the monster simple group, a huge and extraordinary mathematical object, which appears to rely on numerous bizzare coincidences to exists. The axioms for groups give no obvious hint that anything like this exits."}{--- \textup{Richard Borcherds}, Mathematicians: An Outer View...}
\vspace{-20pt}
\epigraph{"The one thing I would really like to know before I die is why the monster group exists."}{--- \textup{John Conway}, in a 2014 interview on Numberphile}
\section{Groups}
\begin{definition}[Binary Operation]
    A binary operation $*$ on a set $G$ is a function $*:G \times G \longrightarrow G$. For any $a, b \in G,$ we will write $a * b$ for $*(a, b)$.
\end{definition}
\begin{eg}
    Modular addition on $\mathbb{Z}/n\mathbb{Z}$, composition of symmetries of equilateral triangle.
\end{eg}
\begin{definition}[Group]
    A group is an ordered pair $(G,*)$, where $G$ is a set and $*$ is a binary operation on $G$ satisfying the following axioms:
    \begin{itemize}
        \item \textbf{Associativity:} $(a * b) * c = a * (b * c);$ for all $a, b, c \in G.$
        \item \textbf{Identity:} There exists an element $e \in G$, called identity of $G$, such that for all $a \in G, a * e = e * a = a.$
        \item \textbf{Inverse:} For all $a \in G$, there exists $b\in G$ such that $a * b = b * a = e.$ We call $b$ is the inverse of $a$ and write $b = a^{-1}.$ 
    \end{itemize}
\end{definition}
\begin{eg}
    $(\mathbb{Q}, +), (\mathbb{Z}, +), (\mathbb{R}, +), (\mathbb{Q}\setminus\{0\}, \times)$.
\end{eg}
\begin{definition}
A group $G$ is called \textit{abelian} (Or \textit{commutative}) if $a * b = b * a,$ for all $a, b \in G$. Otherwise, it is called \textit{non-abelian} (Or \textit{non-commutative}).
\end{definition}
\begin{eg}
    \begin{itemize}
        \item $(R, +)$ is an abelian group.
        \item Group of the symmetries of an equilateral triangle is a non-abelian group.
    \end{itemize}
\end{eg}
Informally, we say $G$ is a group under $*$ if $(G, *)$ is a group, or just $G$ is a group when the operation $*$ is clear from the context.
\begin{eg}
    Consider $G = \{x \in \R: 0 \le x < 1\}.$ \\For all $x, y \in G$ define $*$ via -
    \begin{itemize}
        \item $x*y = x+y - [x+y]$.
        \item $x+y$ (Mod $1)$.
    \end{itemize}
    Prove that $(G, *)$ is an abelian group.
\end{eg}
\begin{explanation}
    For all $x,y \in G,$
    $$
    0 \le x+y \text{ (Mod 1)} < 1 \implies 0 \le x*y < 1\textbf{(Closure)}.
    $$
    For all $x,y,z \in G,$
    \begin{align*}
        (x * y) * z &= (x + y \text{ (Mod 1})) * z\\
                    &= (x + y \text{ (Mod 1})) + z \text{ (Mod 1)}\\
                    &= (x + y) + z \text{ (Mod 1)}\\
                    &= x + (y + z) \text{ (Mod 1)}\\
                    &= x + (y + z \text{ (Mod 1))} \text{ (Mod 1)}\\
                    &= x * (y + z \text{ (Mod 1))}\\
                    &= x * (y * z) \textbf{ (Associativity)}.
    \end{align*}
    For all $x \in G,$
    $$
    x * 0 = x+0 \text{ (Mod 1)} = x.
    $$
    Similarly, $0*x = x$ \textbf{(Identity)}.\\\\
    For all $x \neq0 \in G, (1-x) \in G$ and
    $$
    x * (1-x) = x + 1 - x \text{ Mod 1} = 1 \text{ Mod 1} = 0 \text{ \&}
    $$$$
    (1-x) * x = 1 - x + x \text{ Mod 1} = 1 \text{ Mod 1} = 0
    $$
    And for $0 \in G, 0 * 0 = 0 + 0$ (Mod 1) = 0 \textbf{(Inverse)}\\
    For all $x,y \in G$
    \begin{align*}
        x * y &= x + y \text{ (Mod 1)}\\
              &= y + x \text{ (Mod 1)}\\
              &= y * x
    \end{align*}
    Hence, G is an Abelian group.
\end{explanation}
\section{Matrix Groups}
\begin{eg}
    $GL_n(R) = \{x \in \mathbb{R}^{n \times n}: x \text{ is invertible}\}$ under matrix multiplication forms a group.
\end{eg}
\begin{explanation}
    For all $A, B \in GL_n(R), AB$ is an $n\times n$ matrix. If $A$ and $B$ is invertible, $(AB)^{-1} = B^{-1} A^{-1} \in GL_n(R)$ \textbf{(Closure)}.\\
    Matrix multiplication is associative \textbf{(Associativity)}.\\
    Identity matrix $I_n \in GL_n(R)$ since it is invertible \textbf{(Identity)}.\\
    Inverse exists by defintion \textbf{(Inverse)}.\\
\end{explanation}
\begin{eg}
    $SL_n(R) = \{x \in \mathbb{R}^{n \times n}: det(x) = 1\}$ under matrix multiplication forms a group.
\end{eg}
\begin{explanation}
    For all $A, B \in GL_n(R), AB$ is an $n\times n$ matrix. If det($A$) = det($B$) = 1, then det($AB$) = det($A$) $\times$ det($B$) $= 1 \times 1 = 1$ \textbf{(Closure)}.\\
    Matrix multiplication is associative \textbf{(Associativity)}.\\
    Identity matrix $I_n \in GL_n(R)$ since det($I_n$) = $1$.\\
    For all $A \in GL_n(R)$, inverse exists as det($A$) $\neq 0$ and det($A^{-1}$) $= \frac{1}{\text{det}(A)} = 1$ \textbf{(Inverse)}\\
\end{explanation}